% paper80-copilot-heap-scope-advisors.tex — The Heap and Scope Advisors:
%   Copilot-Guided Heap and Scope Analysis in JuGeo
%   Paper 80 of the Judgment Geometry series.
% Compile: pdflatex paper80-copilot-heap-scope-advisors
\documentclass[11pt]{article}
\usepackage{lmodern}
% jugeo-common.tex — shared preamble for all Judgment Geometry papers
% Include via: % jugeo-common.tex — shared preamble for all Judgment Geometry papers
% Include via: % jugeo-common.tex — shared preamble for all Judgment Geometry papers
% Include via: \input{jugeo-common}

% ── Packages ────────────────────────────────────────────────────────────
\usepackage{amsmath,amssymb,amsthm,mathtools}
\usepackage{tikz-cd}
\usepackage{listings}
\usepackage{xcolor}
\usepackage{xspace}
\usepackage{booktabs}
\usepackage{multirow}
\usepackage{enumitem}
\usepackage[expansion=false]{microtype}
\usepackage{hyperref}
\usepackage{cleveref}
% \usepackage{thmtools}  % disabled: not available in this TeX installation
\usepackage{float}
\usepackage{caption}
\usepackage{subcaption}
\usepackage{tabularx}
\usepackage{stmaryrd}       % \llbracket, \rrbracket
\usepackage{bbm}            % \mathbbm{1}

% ── Page geometry ───────────────────────────────────────────────────────
\usepackage[margin=1in]{geometry}

% ── Theorem environments ────────────────────────────────────────────────
\theoremstyle{plain}
\newtheorem{theorem}{Theorem}[section]
\newtheorem{lemma}[theorem]{Lemma}
\newtheorem{proposition}[theorem]{Proposition}
\newtheorem{corollary}[theorem]{Corollary}

\theoremstyle{definition}
\newtheorem{definition}[theorem]{Definition}
\newtheorem{example}[theorem]{Example}
\newtheorem{construction}[theorem]{Construction}

\theoremstyle{remark}
\newtheorem{remark}[theorem]{Remark}
\newtheorem{notation}[theorem]{Notation}

% ── Python listings style ──────────────────────────────────────────────
\definecolor{jugeoBlue}{HTML}{2563EB}
\definecolor{jugeoGreen}{HTML}{059669}
\definecolor{jugeoRed}{HTML}{DC2626}
\definecolor{jugeoPurple}{HTML}{7C3AED}
\definecolor{jugeoGray}{HTML}{6B7280}
\definecolor{jugeoBg}{HTML}{F8FAFC}

\lstdefinestyle{jugeo-python}{
  language=Python,
  basicstyle=\ttfamily\small,
  keywordstyle=\color{jugeoBlue}\bfseries,
  stringstyle=\color{jugeoGreen},
  commentstyle=\color{jugeoGray}\itshape,
  numberstyle=\tiny\color{jugeoGray},
  backgroundcolor=\color{jugeoBg},
  numbers=left,
  numbersep=6pt,
  frame=single,
  framerule=0.4pt,
  rulecolor=\color{jugeoGray!40},
  breaklines=true,
  breakatwhitespace=true,
  showstringspaces=false,
  tabsize=4,
  xleftmargin=1.5em,
  framexleftmargin=1.5em,
  morekeywords={dataclass,frozen,field,Enum,IntEnum,auto,Any,
                Mapping,Sequence,Optional,match,case},
  literate={→}{{$\to$}}1 {←}{{$\leftarrow$}}1
           {≤}{{$\leq$}}1 {≥}{{$\geq$}}1 {≠}{{$\neq$}}1
           {∈}{{$\in$}}1 {∀}{{$\forall$}}1 {∃}{{$\exists$}}1
           {⊕}{{$\oplus$}}1 {⊖}{{$\ominus$}}1,
  escapeinside={(*@}{@*)},
}
\lstset{style=jugeo-python}

\lstdefinestyle{jugeo-output}{
  basicstyle=\ttfamily\small,
  backgroundcolor=\color{jugeoBg},
  frame=single,
  framerule=0.4pt,
  rulecolor=\color{jugeoGray!40},
  breaklines=true,
  showstringspaces=false,
  xleftmargin=1.5em,
  framexleftmargin=0.5em,
  numbers=none,
}

% ── JuGeo-specific macros ──────────────────────────────────────────────

% Judgment 8-tuple
\newcommand{\J}{\mathcal{J}}
\newcommand{\Jud}[8]{(#1,\,#2,\,#3,\,#4,\,#5,\,#6,\,#7,\,#8)}
\newcommand{\JudShort}[1]{\J_{#1}}

% Components
\newcommand{\coord}{\mathsf{c}}            % coordinate
\newcommand{\prop}{\varphi}                 % proposition
\newcommand{\carrier}{\mathcal{A}}          % carrier type
\newcommand{\evid}{\mathcal{E}}             % evidence bundle
\newcommand{\oblig}{\mathcal{O}}            % obligations
\newcommand{\obstr}{\mathcal{B}}            % obstructions
\newcommand{\trust}{\mathcal{T}}            % trust annotation
\newcommand{\prov}{\Pi}                     % provenance

% Trust algebra
\newcommand{\Talg}{\mathfrak{T}}            % trust ordered algebra
\newcommand{\Eadm}{\mathcal{E}_{\mathrm{adm}}}
\newcommand{\tleq}{\preceq}                 % trust ordering
\newcommand{\tjoin}{\oplus}                 % conservative join
\newcommand{\tatten}{\ominus}               % attenuation
\newcommand{\tpromote}{\uparrow_\pi}        % promotion
\newcommand{\tdemote}{\downarrow_\chi}       % demotion

% Trust tiers
\newcommand{\Tcontra}{\ensuremath{\mathsf{CONTRADICTED}}}
\newcommand{\Tunver}{\ensuremath{\mathsf{UNVERIFIED}}}
\newcommand{\Tcopilot}{\ensuremath{\mathsf{COPILOT}}}
\newcommand{\Toracle}{\ensuremath{\mathsf{ORACLE}}}
\newcommand{\Truntime}{\ensuremath{\mathsf{RUNTIME}}}
\newcommand{\Tsolver}{\ensuremath{\mathsf{SOLVER}}}
\newcommand{\Tproof}{\ensuremath{\mathsf{PROOF}}}

% Site and geometry
\newcommand{\Site}{\mathcal{S}}             % semantic site
\newcommand{\Cov}{\mathrm{Cov}}            % covering family
\newcommand{\Ob}{\mathrm{Ob}}              % objects
\newcommand{\Mor}{\mathrm{Mor}}            % morphisms
\newcommand{\res}{\mathrm{res}}            % restriction
\newcommand{\Cech}{\check{C}}              % Čech complex
\newcommand{\Hcoh}[1]{H^{#1}}             % cohomology group

% Descent
\newcommand{\descent}{\mathrm{desc}}
\newcommand{\glue}{\mathrm{glue}}
\newcommand{\overlap}[2]{#1 \cap #2}

% Semantic moves
\newcommand{\VERIFY}{\textsc{Verify}}
\newcommand{\CONSTRUCT}{\textsc{Construct}}
\newcommand{\NORMALIZE}{\textsc{Normalize}}
\newcommand{\GLUE}{\textsc{Glue}}
\newcommand{\RESTRICT}{\textsc{Restrict}}
\newcommand{\TRANSPORT}{\textsc{Transport}}
\newcommand{\REPAIR}{\textsc{Repair}}
\newcommand{\PROMOTE}{\textsc{Promote}}
\newcommand{\CHALLENGE}{\textsc{Challenge}}
\newcommand{\ESCALATE}{\textsc{Escalate}}

% Fragments
\newcommand{\QFLIA}{\texttt{QF\_LIA}}
\newcommand{\QFLRA}{\texttt{QF\_LRA}}
\newcommand{\QFBV}{\texttt{QF\_BV}}
\newcommand{\QFUF}{\texttt{QF\_UF}}

% Misc
\newcommand{\Python}{\textsc{Python}}
\newcommand{\JuGeo}{\textsc{JuGeo}}
\newcommand{\LEAN}{\textsc{Lean}\,4}
\newcommand{\Fstar}{F$^{\star}$}
\newcommand{\Coq}{\textsc{Coq}}
\newcommand{\Dafny}{\textsc{Dafny}}

% Common shortcuts
\newcommand{\ie}{i.e.\@\xspace}
\newcommand{\eg}{e.g.\@\xspace}
\newcommand{\cf}{cf.\@\xspace}
\newcommand{\wrt}{w.r.t.\@\xspace}

% Evaluation shortcuts
\newcommand{\numcases}{300}
\newcommand{\numspec}{100}
\newcommand{\numeq}{100}
\newcommand{\numbug}{100}
\newcommand{\medtime}{6.5\,ms}
\newcommand{\totaltime}{1.949\,s}


% ── Packages ────────────────────────────────────────────────────────────
\usepackage{amsmath,amssymb,amsthm,mathtools}
\usepackage{tikz-cd}
\usepackage{listings}
\usepackage{xcolor}
\usepackage{xspace}
\usepackage{booktabs}
\usepackage{multirow}
\usepackage{enumitem}
\usepackage[expansion=false]{microtype}
\usepackage{hyperref}
\usepackage{cleveref}
% \usepackage{thmtools}  % disabled: not available in this TeX installation
\usepackage{float}
\usepackage{caption}
\usepackage{subcaption}
\usepackage{tabularx}
\usepackage{stmaryrd}       % \llbracket, \rrbracket
\usepackage{bbm}            % \mathbbm{1}

% ── Page geometry ───────────────────────────────────────────────────────
\usepackage[margin=1in]{geometry}

% ── Theorem environments ────────────────────────────────────────────────
\theoremstyle{plain}
\newtheorem{theorem}{Theorem}[section]
\newtheorem{lemma}[theorem]{Lemma}
\newtheorem{proposition}[theorem]{Proposition}
\newtheorem{corollary}[theorem]{Corollary}

\theoremstyle{definition}
\newtheorem{definition}[theorem]{Definition}
\newtheorem{example}[theorem]{Example}
\newtheorem{construction}[theorem]{Construction}

\theoremstyle{remark}
\newtheorem{remark}[theorem]{Remark}
\newtheorem{notation}[theorem]{Notation}

% ── Python listings style ──────────────────────────────────────────────
\definecolor{jugeoBlue}{HTML}{2563EB}
\definecolor{jugeoGreen}{HTML}{059669}
\definecolor{jugeoRed}{HTML}{DC2626}
\definecolor{jugeoPurple}{HTML}{7C3AED}
\definecolor{jugeoGray}{HTML}{6B7280}
\definecolor{jugeoBg}{HTML}{F8FAFC}

\lstdefinestyle{jugeo-python}{
  language=Python,
  basicstyle=\ttfamily\small,
  keywordstyle=\color{jugeoBlue}\bfseries,
  stringstyle=\color{jugeoGreen},
  commentstyle=\color{jugeoGray}\itshape,
  numberstyle=\tiny\color{jugeoGray},
  backgroundcolor=\color{jugeoBg},
  numbers=left,
  numbersep=6pt,
  frame=single,
  framerule=0.4pt,
  rulecolor=\color{jugeoGray!40},
  breaklines=true,
  breakatwhitespace=true,
  showstringspaces=false,
  tabsize=4,
  xleftmargin=1.5em,
  framexleftmargin=1.5em,
  morekeywords={dataclass,frozen,field,Enum,IntEnum,auto,Any,
                Mapping,Sequence,Optional,match,case},
  literate={→}{{$\to$}}1 {←}{{$\leftarrow$}}1
           {≤}{{$\leq$}}1 {≥}{{$\geq$}}1 {≠}{{$\neq$}}1
           {∈}{{$\in$}}1 {∀}{{$\forall$}}1 {∃}{{$\exists$}}1
           {⊕}{{$\oplus$}}1 {⊖}{{$\ominus$}}1,
  escapeinside={(*@}{@*)},
}
\lstset{style=jugeo-python}

\lstdefinestyle{jugeo-output}{
  basicstyle=\ttfamily\small,
  backgroundcolor=\color{jugeoBg},
  frame=single,
  framerule=0.4pt,
  rulecolor=\color{jugeoGray!40},
  breaklines=true,
  showstringspaces=false,
  xleftmargin=1.5em,
  framexleftmargin=0.5em,
  numbers=none,
}

% ── JuGeo-specific macros ──────────────────────────────────────────────

% Judgment 8-tuple
\newcommand{\J}{\mathcal{J}}
\newcommand{\Jud}[8]{(#1,\,#2,\,#3,\,#4,\,#5,\,#6,\,#7,\,#8)}
\newcommand{\JudShort}[1]{\J_{#1}}

% Components
\newcommand{\coord}{\mathsf{c}}            % coordinate
\newcommand{\prop}{\varphi}                 % proposition
\newcommand{\carrier}{\mathcal{A}}          % carrier type
\newcommand{\evid}{\mathcal{E}}             % evidence bundle
\newcommand{\oblig}{\mathcal{O}}            % obligations
\newcommand{\obstr}{\mathcal{B}}            % obstructions
\newcommand{\trust}{\mathcal{T}}            % trust annotation
\newcommand{\prov}{\Pi}                     % provenance

% Trust algebra
\newcommand{\Talg}{\mathfrak{T}}            % trust ordered algebra
\newcommand{\Eadm}{\mathcal{E}_{\mathrm{adm}}}
\newcommand{\tleq}{\preceq}                 % trust ordering
\newcommand{\tjoin}{\oplus}                 % conservative join
\newcommand{\tatten}{\ominus}               % attenuation
\newcommand{\tpromote}{\uparrow_\pi}        % promotion
\newcommand{\tdemote}{\downarrow_\chi}       % demotion

% Trust tiers
\newcommand{\Tcontra}{\ensuremath{\mathsf{CONTRADICTED}}}
\newcommand{\Tunver}{\ensuremath{\mathsf{UNVERIFIED}}}
\newcommand{\Tcopilot}{\ensuremath{\mathsf{COPILOT}}}
\newcommand{\Toracle}{\ensuremath{\mathsf{ORACLE}}}
\newcommand{\Truntime}{\ensuremath{\mathsf{RUNTIME}}}
\newcommand{\Tsolver}{\ensuremath{\mathsf{SOLVER}}}
\newcommand{\Tproof}{\ensuremath{\mathsf{PROOF}}}

% Site and geometry
\newcommand{\Site}{\mathcal{S}}             % semantic site
\newcommand{\Cov}{\mathrm{Cov}}            % covering family
\newcommand{\Ob}{\mathrm{Ob}}              % objects
\newcommand{\Mor}{\mathrm{Mor}}            % morphisms
\newcommand{\res}{\mathrm{res}}            % restriction
\newcommand{\Cech}{\check{C}}              % Čech complex
\newcommand{\Hcoh}[1]{H^{#1}}             % cohomology group

% Descent
\newcommand{\descent}{\mathrm{desc}}
\newcommand{\glue}{\mathrm{glue}}
\newcommand{\overlap}[2]{#1 \cap #2}

% Semantic moves
\newcommand{\VERIFY}{\textsc{Verify}}
\newcommand{\CONSTRUCT}{\textsc{Construct}}
\newcommand{\NORMALIZE}{\textsc{Normalize}}
\newcommand{\GLUE}{\textsc{Glue}}
\newcommand{\RESTRICT}{\textsc{Restrict}}
\newcommand{\TRANSPORT}{\textsc{Transport}}
\newcommand{\REPAIR}{\textsc{Repair}}
\newcommand{\PROMOTE}{\textsc{Promote}}
\newcommand{\CHALLENGE}{\textsc{Challenge}}
\newcommand{\ESCALATE}{\textsc{Escalate}}

% Fragments
\newcommand{\QFLIA}{\texttt{QF\_LIA}}
\newcommand{\QFLRA}{\texttt{QF\_LRA}}
\newcommand{\QFBV}{\texttt{QF\_BV}}
\newcommand{\QFUF}{\texttt{QF\_UF}}

% Misc
\newcommand{\Python}{\textsc{Python}}
\newcommand{\JuGeo}{\textsc{JuGeo}}
\newcommand{\LEAN}{\textsc{Lean}\,4}
\newcommand{\Fstar}{F$^{\star}$}
\newcommand{\Coq}{\textsc{Coq}}
\newcommand{\Dafny}{\textsc{Dafny}}

% Common shortcuts
\newcommand{\ie}{i.e.\@\xspace}
\newcommand{\eg}{e.g.\@\xspace}
\newcommand{\cf}{cf.\@\xspace}
\newcommand{\wrt}{w.r.t.\@\xspace}

% Evaluation shortcuts
\newcommand{\numcases}{300}
\newcommand{\numspec}{100}
\newcommand{\numeq}{100}
\newcommand{\numbug}{100}
\newcommand{\medtime}{6.5\,ms}
\newcommand{\totaltime}{1.949\,s}


% ── Packages ────────────────────────────────────────────────────────────
\usepackage{amsmath,amssymb,amsthm,mathtools}
\usepackage{tikz-cd}
\usepackage{listings}
\usepackage{xcolor}
\usepackage{xspace}
\usepackage{booktabs}
\usepackage{multirow}
\usepackage{enumitem}
\usepackage[expansion=false]{microtype}
\usepackage{hyperref}
\usepackage{cleveref}
% \usepackage{thmtools}  % disabled: not available in this TeX installation
\usepackage{float}
\usepackage{caption}
\usepackage{subcaption}
\usepackage{tabularx}
\usepackage{stmaryrd}       % \llbracket, \rrbracket
\usepackage{bbm}            % \mathbbm{1}

% ── Page geometry ───────────────────────────────────────────────────────
\usepackage[margin=1in]{geometry}

% ── Theorem environments ────────────────────────────────────────────────
\theoremstyle{plain}
\newtheorem{theorem}{Theorem}[section]
\newtheorem{lemma}[theorem]{Lemma}
\newtheorem{proposition}[theorem]{Proposition}
\newtheorem{corollary}[theorem]{Corollary}

\theoremstyle{definition}
\newtheorem{definition}[theorem]{Definition}
\newtheorem{example}[theorem]{Example}
\newtheorem{construction}[theorem]{Construction}

\theoremstyle{remark}
\newtheorem{remark}[theorem]{Remark}
\newtheorem{notation}[theorem]{Notation}

% ── Python listings style ──────────────────────────────────────────────
\definecolor{jugeoBlue}{HTML}{2563EB}
\definecolor{jugeoGreen}{HTML}{059669}
\definecolor{jugeoRed}{HTML}{DC2626}
\definecolor{jugeoPurple}{HTML}{7C3AED}
\definecolor{jugeoGray}{HTML}{6B7280}
\definecolor{jugeoBg}{HTML}{F8FAFC}

\lstdefinestyle{jugeo-python}{
  language=Python,
  basicstyle=\ttfamily\small,
  keywordstyle=\color{jugeoBlue}\bfseries,
  stringstyle=\color{jugeoGreen},
  commentstyle=\color{jugeoGray}\itshape,
  numberstyle=\tiny\color{jugeoGray},
  backgroundcolor=\color{jugeoBg},
  numbers=left,
  numbersep=6pt,
  frame=single,
  framerule=0.4pt,
  rulecolor=\color{jugeoGray!40},
  breaklines=true,
  breakatwhitespace=true,
  showstringspaces=false,
  tabsize=4,
  xleftmargin=1.5em,
  framexleftmargin=1.5em,
  morekeywords={dataclass,frozen,field,Enum,IntEnum,auto,Any,
                Mapping,Sequence,Optional,match,case},
  literate={→}{{$\to$}}1 {←}{{$\leftarrow$}}1
           {≤}{{$\leq$}}1 {≥}{{$\geq$}}1 {≠}{{$\neq$}}1
           {∈}{{$\in$}}1 {∀}{{$\forall$}}1 {∃}{{$\exists$}}1
           {⊕}{{$\oplus$}}1 {⊖}{{$\ominus$}}1,
  escapeinside={(*@}{@*)},
}
\lstset{style=jugeo-python}

\lstdefinestyle{jugeo-output}{
  basicstyle=\ttfamily\small,
  backgroundcolor=\color{jugeoBg},
  frame=single,
  framerule=0.4pt,
  rulecolor=\color{jugeoGray!40},
  breaklines=true,
  showstringspaces=false,
  xleftmargin=1.5em,
  framexleftmargin=0.5em,
  numbers=none,
}

% ── JuGeo-specific macros ──────────────────────────────────────────────

% Judgment 8-tuple
\newcommand{\J}{\mathcal{J}}
\newcommand{\Jud}[8]{(#1,\,#2,\,#3,\,#4,\,#5,\,#6,\,#7,\,#8)}
\newcommand{\JudShort}[1]{\J_{#1}}

% Components
\newcommand{\coord}{\mathsf{c}}            % coordinate
\newcommand{\prop}{\varphi}                 % proposition
\newcommand{\carrier}{\mathcal{A}}          % carrier type
\newcommand{\evid}{\mathcal{E}}             % evidence bundle
\newcommand{\oblig}{\mathcal{O}}            % obligations
\newcommand{\obstr}{\mathcal{B}}            % obstructions
\newcommand{\trust}{\mathcal{T}}            % trust annotation
\newcommand{\prov}{\Pi}                     % provenance

% Trust algebra
\newcommand{\Talg}{\mathfrak{T}}            % trust ordered algebra
\newcommand{\Eadm}{\mathcal{E}_{\mathrm{adm}}}
\newcommand{\tleq}{\preceq}                 % trust ordering
\newcommand{\tjoin}{\oplus}                 % conservative join
\newcommand{\tatten}{\ominus}               % attenuation
\newcommand{\tpromote}{\uparrow_\pi}        % promotion
\newcommand{\tdemote}{\downarrow_\chi}       % demotion

% Trust tiers
\newcommand{\Tcontra}{\ensuremath{\mathsf{CONTRADICTED}}}
\newcommand{\Tunver}{\ensuremath{\mathsf{UNVERIFIED}}}
\newcommand{\Tcopilot}{\ensuremath{\mathsf{COPILOT}}}
\newcommand{\Toracle}{\ensuremath{\mathsf{ORACLE}}}
\newcommand{\Truntime}{\ensuremath{\mathsf{RUNTIME}}}
\newcommand{\Tsolver}{\ensuremath{\mathsf{SOLVER}}}
\newcommand{\Tproof}{\ensuremath{\mathsf{PROOF}}}

% Site and geometry
\newcommand{\Site}{\mathcal{S}}             % semantic site
\newcommand{\Cov}{\mathrm{Cov}}            % covering family
\newcommand{\Ob}{\mathrm{Ob}}              % objects
\newcommand{\Mor}{\mathrm{Mor}}            % morphisms
\newcommand{\res}{\mathrm{res}}            % restriction
\newcommand{\Cech}{\check{C}}              % Čech complex
\newcommand{\Hcoh}[1]{H^{#1}}             % cohomology group

% Descent
\newcommand{\descent}{\mathrm{desc}}
\newcommand{\glue}{\mathrm{glue}}
\newcommand{\overlap}[2]{#1 \cap #2}

% Semantic moves
\newcommand{\VERIFY}{\textsc{Verify}}
\newcommand{\CONSTRUCT}{\textsc{Construct}}
\newcommand{\NORMALIZE}{\textsc{Normalize}}
\newcommand{\GLUE}{\textsc{Glue}}
\newcommand{\RESTRICT}{\textsc{Restrict}}
\newcommand{\TRANSPORT}{\textsc{Transport}}
\newcommand{\REPAIR}{\textsc{Repair}}
\newcommand{\PROMOTE}{\textsc{Promote}}
\newcommand{\CHALLENGE}{\textsc{Challenge}}
\newcommand{\ESCALATE}{\textsc{Escalate}}

% Fragments
\newcommand{\QFLIA}{\texttt{QF\_LIA}}
\newcommand{\QFLRA}{\texttt{QF\_LRA}}
\newcommand{\QFBV}{\texttt{QF\_BV}}
\newcommand{\QFUF}{\texttt{QF\_UF}}

% Misc
\newcommand{\Python}{\textsc{Python}}
\newcommand{\JuGeo}{\textsc{JuGeo}}
\newcommand{\LEAN}{\textsc{Lean}\,4}
\newcommand{\Fstar}{F$^{\star}$}
\newcommand{\Coq}{\textsc{Coq}}
\newcommand{\Dafny}{\textsc{Dafny}}

% Common shortcuts
\newcommand{\ie}{i.e.\@\xspace}
\newcommand{\eg}{e.g.\@\xspace}
\newcommand{\cf}{cf.\@\xspace}
\newcommand{\wrt}{w.r.t.\@\xspace}

% Evaluation shortcuts
\newcommand{\numcases}{300}
\newcommand{\numspec}{100}
\newcommand{\numeq}{100}
\newcommand{\numbug}{100}
\newcommand{\medtime}{6.5\,ms}
\newcommand{\totaltime}{1.949\,s}


% ── Additional packages ────────────────────────────────────────────
\usepackage{array}
\usepackage{tikz}
\usetikzlibrary{arrows.meta,positioning,calc,fit,backgrounds}

% ── Paper-specific macros ──────────────────────────────────────────
\newcommand{\HeapAdv}{\mathsf{CopilotHeapAdvisor}}
\newcommand{\ScopeAdv}{\mathsf{CopilotScopeAdvisor}}
\newcommand{\AdvCompose}{\mathsf{AdvisorComposition}}
\newcommand{\HeapPart}{\mathsf{HeapPartition}}
\newcommand{\ScopeRes}{\mathsf{ScopeResolution}}
\newcommand{\HeapQuery}{\mathrm{heapQuery}}
\newcommand{\ScopeQuery}{\mathrm{scopeQuery}}
\newcommand{\AdvSound}{\mathrm{sound}}
\newcommand{\AdvComplete}{\mathrm{complete}}
\newcommand{\AdvTrust}{\mathrm{advisorTrust}}
\newcommand{\CopChan}{\mathsf{CopilotChannel}}
\newcommand{\CopBridge}{\mathsf{CopilotAPIBridge}}
\newcommand{\TrustCeil}{\mathrm{ceiling}}
\newcommand{\HeapRegion}{\mathsf{HeapRegion}}
\newcommand{\ScopeChain}{\mathsf{ScopeChain}}
\newcommand{\MutAnalysis}{\mathrm{mutationAnalysis}}
\newcommand{\ShadowDetect}{\mathrm{shadowDetect}}

\title{\textbf{The Heap and Scope Advisors:\\
Copilot-Guided Heap and Scope Analysis in Judgment Geometry}}

\author{Halley Young}

\date{Paper~80 --- Judgment Geometry Series}

\begin{document}
\maketitle

% ─────────────────────────────────────────────────────────────────────
\begin{abstract}
Heap aliasing and scope shadowing are pervasive sources of bugs in
imperative programs, yet their analysis often requires domain-specific
insight that generic static analyses cannot provide.  We introduce
two Copilot-guided advisors within the \JuGeo{} framework:
(i)~the $\HeapAdv$, which uses LLM inference to suggest heap
partitions, identify aliasing patterns, and flag mutation-related
anomalies; and
(ii)~the $\ScopeAdv$, which uses LLM inference to detect scope
shadowing, identify variable capture issues, and suggest scope
restructurings.

Both advisors operate within the trust-calibrated evidence
architecture of \JuGeo{}, producing advice at the $\Tcopilot$
trust tier.  All advice must be independently validated before
influencing the judgment state.  We formalize the notion of
advisory soundness (advice that passes validation is correct),
prove that advisory composition preserves trust bounds, and
establish that the advisor pattern separates concern identification
from concern verification.

We prove five main theorems: heap advisory soundness, scope advisory
soundness, composition trust preservation, advisory completeness
relative to a ground-truth oracle, and health monitoring correctness
for advisor subsystems.  All proofs are presented in full and
formalized in \LEAN.
\end{abstract}

% ─────────────────────────────────────────────────────────────────────
\section{Introduction}
\label{sec:intro}

Heap aliasing and scope shadowing represent two of the most
common classes of subtle bugs in imperative and object-oriented
programs.  Heap aliasing occurs when multiple references point to
the same memory location, leading to unexpected mutations.
Scope shadowing occurs when a variable in an inner scope hides
a variable in an outer scope, leading to unintended name
resolution.

Static analysis tools can detect some instances of these
problems~\cite{hoare1969axiomatic}, but many cases require
understanding the programmer's \emph{intent}---information that
is not captured in the program text alone.  Large language models,
having been trained on vast codebases with associated documentation
and commit messages, can infer intent to a degree that purely
syntactic analyses cannot.

In the \JuGeo{} framework, \emph{advisors} are components that
generate domain-specific guidance for the verification process.
Each advisor uses the $\CopChan$ to query the LLM and produces
advice that is calibrated to the $\Tcopilot$ trust tier.  The
advice is then validated by the solver or runtime channel before
it can influence the judgment state.

\paragraph{Contributions.}
\begin{itemize}[leftmargin=2em]
\item We define the $\HeapAdv$ for LLM-guided heap analysis,
  including heap partitioning, aliasing detection, and mutation
  flagging (\S\ref{sec:heap}).

\item We define the $\ScopeAdv$ for LLM-guided scope analysis,
  including shadowing detection, capture identification, and
  scope restructuring (\S\ref{sec:scope}).

\item We formalize the advisory pattern and prove soundness,
  trust preservation, and completeness properties
  (\S\ref{sec:formal}).

\item We show how both advisors integrate with the \JuGeo{}
  health monitoring infrastructure (\S\ref{sec:health}).

\item We describe the composition of multiple advisors and prove
  that composition preserves trust bounds (\S\ref{sec:composition}).
\end{itemize}

\paragraph{Outline.}
Section~\ref{sec:pattern} describes the advisor pattern.
Section~\ref{sec:heap} presents the heap advisor.
Section~\ref{sec:scope} presents the scope advisor.
Section~\ref{sec:composition} covers advisor composition.
Section~\ref{sec:health} describes health monitoring.
Section~\ref{sec:formal} proves the main theorems.
Section~\ref{sec:discussion} discusses implications.
Section~\ref{sec:related} surveys related work.
Section~\ref{sec:conclusion} concludes.

% ─────────────────────────────────────────────────────────────────────
\section{The Advisor Pattern}
\label{sec:pattern}

\subsection{Advisory Interface}
\label{sec:pattern:interface}

\begin{definition}[Advisor]
\label{def:advisor}
An \emph{advisor} $A$ is a component that:
\begin{enumerate}
\item Accepts a verification context $(c, \varphi, \Gamma)$
  consisting of a coordinate $c$, proposition $\varphi$, and
  environment $\Gamma$.
\item Queries the $\CopChan$ with a domain-specific prompt
  constructed from the context.
\item Produces a set of \emph{advice items} $\{a_1, \ldots, a_k\}$,
  each carrying trust $\Tcopilot$.
\item Returns the advice items for validation by the verification
  engine.
\end{enumerate}
\end{definition}

\begin{definition}[Advice item]
\label{def:advice-item}
An \emph{advice item} is a tuple $a = (d, s, \gamma, \tau)$ where:
\begin{itemize}[leftmargin=2em]
\item $d$ is the domain (e.g., ``heap'', ``scope'').
\item $s$ is the suggested action (a structured description).
\item $\gamma \in [0, 1]$ is a confidence score.
\item $\tau = \Tcopilot$ is the trust level.
\end{itemize}
\end{definition}

\begin{definition}[Advisory soundness]
\label{def:adv-sound}
An advisor $A$ is \emph{sound} if for every advice item $a$
produced by $A$ that passes validation:
$\DeductValid(a) \implies \mathrm{correct}(a)$,
\ie validated advice items correspond to genuine issues or
valid suggestions.
\end{definition}

\subsection{Advisor Trust Model}
\label{sec:pattern:trust}

\begin{lemma}[Advisor trust ceiling]
\label{lem:adv-trust}
Every advice item produced by an advisor using the $\CopChan$
has trust at most $\Tcopilot$:
$\forall a \in A(c, \varphi, \Gamma).\; a.\tau \tleq \Tcopilot$.
\end{lemma}

\begin{proof}
The advisor queries the $\CopChan$, which enforces the trust
ceiling $\Tcopilot$ on all outputs.  The advisor does not
modify trust levels.  Therefore all advice items inherit the
channel's trust ceiling.
\end{proof}

% ─────────────────────────────────────────────────────────────────────
\section{The Copilot Heap Advisor}
\label{sec:heap}

\subsection{Heap Analysis Tasks}
\label{sec:heap:tasks}

The $\HeapAdv$ addresses three categories of heap-related issues:

\begin{definition}[Heap partition]
\label{def:heap-part}
A \emph{heap partition} is a decomposition of the heap into
disjoint regions $R_1, R_2, \ldots, R_m$ such that:
\begin{enumerate}
\item Each region contains a set of related objects.
\item Cross-region aliasing is flagged as a potential issue.
\item Mutations within a region are tracked for consistency.
\end{enumerate}
\end{definition}

\begin{definition}[Aliasing pattern]
\label{def:aliasing}
An \emph{aliasing pattern} is a pair $(r_1, r_2)$ of references
where $r_1$ and $r_2$ may point to the same heap object.  An
aliasing pattern is \emph{benign} if the program's correctness
does not depend on aliasing status, and \emph{hazardous} if
mutation through one reference may invalidate assumptions about
the other.
\end{definition}

\begin{definition}[Mutation anomaly]
\label{def:mutation}
A \emph{mutation anomaly} is a heap modification that violates
an expected invariant, such as:
\begin{itemize}[leftmargin=2em]
\item Mutation of a logically immutable object.
\item Mutation through an aliased reference without synchronization.
\item Mutation after a reference has been shared with another scope.
\end{itemize}
\end{definition}

The following listing shows how the $\HeapAdv$ uses the
$\CopChan$ for heap analysis.

\begin{lstlisting}[style=jugeo-python, caption={CopilotChannel for heap analysis advisory.}, label=lst:heap-query]
from jugeo.evidence.channels import (
    CopilotChannel,
    ChannelConfiguration,
    EvidenceChannel,
    ChannelJurisdiction,
    EvidenceRequest,
    EvidenceResponse,
)
from jugeo.interfaces.api import (
    CopilotAPIBridge,
    APIEventLog,
    JuGeoAPI,
    APIRequest,
    OperationKind,
)

# Configure Copilot channel for heap analysis
jurisdiction = ChannelJurisdiction.for_channel(EvidenceChannel.COPILOT)
config = ChannelConfiguration(
    channel=EvidenceChannel.COPILOT,
    timeout_ms=12000,
    max_retries=2,
    batch_size=1,
    rate_limit=8.0,
    trust_ceiling="proposal",
    jurisdiction=jurisdiction,
    is_enabled=True,
    priority=35,
)
copilot = CopilotChannel(config=config, model_name="copilot-default")

# Query for heap partition suggestions
heap_prompt = (
    "Analyze the following code for heap aliasing issues:\n"
    "  def process(data, cache):\n"
    "      result = cache.get(data.key)\n"
    "      if result is None:\n"
    "          result = compute(data)\n"
    "          cache[data.key] = result\n"
    "      data.result = result  # potential aliasing\n"
    "      return data\n\n"
    "Identify: (1) heap regions, (2) aliasing patterns,\n"
    "(3) mutation anomalies."
)
heap_result = copilot.query_llm(prompt=heap_prompt, timeout_ms=10000)

# Build evidence request for heap analysis
heap_request = EvidenceRequest(
    request_id="heap-advisory-001",
    coordinate="process.cache.aliasing",
    proposition="cache[key] and data.result do not alias",
    required_kind="PROPOSAL",
    proposition_kind="heap_analysis",
    preferred_channel=EvidenceChannel.COPILOT,
    fallback_channels=(EvidenceChannel.RUNTIME,),
    deadline_ms=12000.0,
    budget=1.5,
    metadata={"advisor": "heap", "analysis_type": "aliasing"},
)
heap_response = copilot.handle_request(heap_request)

# Verify trust ceiling is respected
assert heap_response.trust_level == "proposal"
print(f"Heap advisory trust: {heap_response.trust_level}")
print(f"Advisory is partial: {heap_response.is_partial}")
print(f"Residuals: {heap_response.residuals}")

# Set up API bridge for heap advisory pipeline
event_log = APIEventLog()
bridge = CopilotAPIBridge(
    channel=copilot,
    event_log=event_log,
    require_corroboration=True,
)
api = JuGeoAPI(event_log=event_log, copilot_bridge=bridge)

# Use the API for inspection-based heap analysis
inspect_request = APIRequest(
    request_id="heap-inspect-001",
    operation=OperationKind.INSPECT,
    coordinate="process.cache",
    proposition="no hazardous aliasing in cache operations",
    budget=5000,
    deadline=15.0,
    metadata={"advisor": "heap"},
)
inspect_response = api.inspect(inspect_request)
print(f"Inspection status: {inspect_response.status}")
print(f"Copilot contributed: {inspect_response.copilot_contributed}")
\end{lstlisting}

\subsection{Heap Advisory Workflow}
\label{sec:heap:workflow}

The $\HeapAdv$ workflow proceeds in three phases:

\begin{enumerate}
\item \textbf{Partition suggestion.}  The LLM analyzes the code
  and suggests a heap partition.  Each suggested region includes
  a description and a set of references.

\item \textbf{Aliasing detection.}  For each pair of regions,
  the LLM identifies potential cross-region aliases.  Each alias
  is classified as benign or hazardous.

\item \textbf{Mutation flagging.}  The LLM reviews mutation
  operations and flags those that may violate heap invariants.
  Each flag includes the mutation site, the violated invariant,
  and a suggested fix.
\end{enumerate}

\begin{remark}[Heap advisor scope]
\label{rem:heap-scope}
The $\HeapAdv$ does not perform the actual heap analysis; it
\emph{suggests} where to look and what patterns to check.  The
actual verification is delegated to the solver or runtime channel.
This separation ensures that the advisor's output (at $\Tcopilot$)
does not directly affect the judgment state.
\end{remark}

% ─────────────────────────────────────────────────────────────────────
\section{The Copilot Scope Advisor}
\label{sec:scope}

\subsection{Scope Analysis Tasks}
\label{sec:scope:tasks}

The $\ScopeAdv$ addresses three categories of scope-related issues:

\begin{definition}[Scope shadowing]
\label{def:shadowing}
\emph{Scope shadowing} occurs when a variable $x$ declared in
scope $S_2$ has the same name as a variable in enclosing scope
$S_1$, causing the inner declaration to hide the outer one.
Shadowing is \emph{intentional} if the programmer deliberately
chose to reuse the name, and \emph{accidental} if it results
from oversight.
\end{definition}

\begin{definition}[Variable capture]
\label{def:capture}
\emph{Variable capture} occurs when a closure or lambda expression
captures a reference to a mutable variable in an enclosing scope.
Subsequent mutations of the variable affect the closure's behavior,
potentially violating its expected semantics.
\end{definition}

\begin{definition}[Scope restructuring]
\label{def:restructuring}
A \emph{scope restructuring} is a transformation of the program's
scope structure that eliminates shadowing or capture issues without
changing the program's semantics.
\end{definition}

The following listing shows how the $\ScopeAdv$ queries the
$\CopChan$ and integrates with health monitoring.

\begin{lstlisting}[style=jugeo-python, caption={Scope advisor with CopilotAPIBridge and health monitoring.}, label=lst:scope-query]
from jugeo.evidence.channels import (
    CopilotChannel,
    ChannelConfiguration,
    EvidenceChannel,
    EvidenceRequest,
    EvidenceResponse,
    ChannelMonitor,
)
from jugeo.interfaces.api import (
    CopilotAPIBridge,
    APIEventLog,
    JuGeoAPI,
    APIRequest,
    APIResponse,
    OperationKind,
    RequestStatus,
)
from jugeo.kernel.health import (
    HealthDimension,
    HealthIndicator,
    HealthStatus,
    HealthReport,
)

# Configure Copilot for scope analysis
config = ChannelConfiguration.default_for(EvidenceChannel.COPILOT)
copilot = CopilotChannel(config=config, model_name="copilot-default")
monitor = ChannelMonitor()

# Query for scope shadowing detection
scope_prompt = (
    "Analyze scope shadowing in this Python function:\n"
    "  x = 10\n"
    "  def outer():\n"
    "      x = 20  # shadows module-level x\n"
    "      def inner():\n"
    "          x = 30  # shadows outer x\n"
    "          return x\n"
    "      return inner() + x\n"
    "Identify all shadowing instances and classify as\n"
    "intentional or accidental."
)
scope_result = copilot.query_llm(prompt=scope_prompt, timeout_ms=8000)

# Build scope advisory request
scope_request = EvidenceRequest(
    request_id="scope-advisory-001",
    coordinate="module.outer.inner.shadowing",
    proposition="variable x is not accidentally shadowed",
    required_kind="PROPOSAL",
    proposition_kind="scope_analysis",
    preferred_channel=EvidenceChannel.COPILOT,
    fallback_channels=(),
    deadline_ms=8000.0,
    budget=1.0,
    metadata={"advisor": "scope", "issue": "shadowing"},
)

monitor.record_request(scope_request)
scope_response = copilot.handle_request(scope_request)
monitor.record_response(scope_response)

# Check advisor health via health indicators
copilot_health = HealthIndicator(
    dimension=HealthDimension.COPILOT_CONNECTIVITY,
    status=HealthStatus.HEALTHY,
    message="Copilot scope advisor responding normally",
    value=1.0,
    timestamp=0.0,
    details={"advisor": "scope", "latency_ms": scope_response.latency_ms},
)
print(f"Health indicator: {copilot_health.dimension}")
print(f"Status: {copilot_health.status}")
print(f"Is healthy: {copilot_health.is_healthy()}")
print(f"Is actionable: {copilot_health.is_actionable()}")

# Build health report for the scope advisor subsystem
report = HealthReport(
    overall_status=HealthStatus.HEALTHY,
    indicators=(copilot_health,),
    generated_at=0.0,
    report_id="scope-advisor-health-001",
    degradation_reasons=(),
    recovery_suggestions=(),
)
print(f"Report healthy: {report.is_fully_healthy()}")
print(f"Worst indicator: {report.worst_indicator}")

# Monitor performance
print(f"Success rate: {monitor.success_rate()}")
print(f"Latency: {monitor.latency_percentiles()}")

# Set up full API pipeline for scope advisory
event_log = APIEventLog()
bridge = CopilotAPIBridge(
    channel=copilot,
    event_log=event_log,
    require_corroboration=True,
)
api = JuGeoAPI(event_log=event_log, copilot_bridge=bridge)

# Verify scope safety through the API
verify_request = APIRequest(
    request_id="scope-verify-001",
    operation=OperationKind.VERIFY,
    coordinate="module.outer.inner",
    proposition="no accidental variable shadowing",
    budget=10000,
    deadline=20.0,
    metadata={"advisor": "scope"},
)
verify_response = api.verify(verify_request)
print(f"Verification: {verify_response.status}")
\end{lstlisting}

\subsection{Scope Advisory Workflow}
\label{sec:scope:workflow}

The $\ScopeAdv$ workflow:

\begin{enumerate}
\item \textbf{Shadowing scan.}  The LLM scans the code for all
  instances of scope shadowing, classifying each as intentional
  or accidental.

\item \textbf{Capture analysis.}  For each closure or lambda,
  the LLM identifies captured variables and flags mutable captures.

\item \textbf{Restructuring suggestion.}  When shadowing or capture
  issues are detected, the LLM suggests scope restructurings that
  preserve semantics while eliminating the issues.
\end{enumerate}

% ─────────────────────────────────────────────────────────────────────
\section{Advisor Composition}
\label{sec:composition}

\subsection{Composing Heap and Scope Advisors}
\label{sec:composition:combine}

\begin{definition}[Advisor composition]
\label{def:adv-compose}
The \emph{composition} of advisors $A_1, A_2$ is the advisor
$A_1 \otimes A_2$ that:
\begin{enumerate}
\item Runs both $A_1$ and $A_2$ on the same context.
\item Merges the resulting advice sets: $\{a_1, \ldots, a_j\}
  \cup \{b_1, \ldots, b_k\}$.
\item Assigns the composed trust as the meet:
  $\tau(A_1 \otimes A_2) = \tau(A_1) \tjoin \tau(A_2)$.
\end{enumerate}
\end{definition}

\begin{lemma}[Composition preserves trust ceiling]
\label{lem:comp-trust}
For any advisors $A_1, A_2$ using the $\CopChan$:
$\tau(A_1 \otimes A_2) = \Tcopilot$.
\end{lemma}

\begin{proof}
By Lemma~\ref{lem:adv-trust}, $\tau(A_1) = \tau(A_2) = \Tcopilot$.
The meet is $\Tcopilot \tjoin \Tcopilot = \Tcopilot$.
\end{proof}

\begin{definition}[Cross-advisor interaction]
\label{def:cross-advisor}
A \emph{cross-advisor interaction} occurs when an advice item
from $A_1$ relates to an advice item from $A_2$.  For the
$\HeapAdv$ and $\ScopeAdv$, interactions arise when:
\begin{itemize}[leftmargin=2em]
\item A heap aliasing issue involves references from different scopes.
\item A scope shadowing affects the resolution of a heap reference.
\item A mutation anomaly occurs in a shadowed variable.
\end{itemize}
\end{definition}

\begin{remark}[Independent analysis]
\label{rem:independent}
In the current architecture, advisors analyze independently.
Cross-advisor interactions are detected during the validation
phase, where the solver or runtime channel may identify that
two advice items from different advisors address related issues.
A future extension could support correlated multi-advisor analysis.
\end{remark}

% ─────────────────────────────────────────────────────────────────────
\section{Health Monitoring for Advisors}
\label{sec:health}

\subsection{Advisor Health Dimensions}
\label{sec:health:dimensions}

The \JuGeo{} health monitoring infrastructure tracks advisor
performance across several dimensions:

\begin{definition}[Advisor health model]
\label{def:advisor-health}
The health of an advisor is tracked along these dimensions:
\begin{itemize}[leftmargin=2em]
\item $\mathsf{COPILOT\_CONNECTIVITY}$: whether the underlying
  $\CopChan$ is responding within timeout bounds.
\item $\mathsf{EVIDENCE\_FLOW}$: whether the advisor is producing
  evidence at a reasonable rate.
\item $\mathsf{TRUST\_CONSISTENCY}$: whether produced evidence
  respects the trust ceiling.
\end{itemize}
\end{definition}

The following listing shows how health indicators are constructed
and monitored for the advisor subsystem.

\begin{lstlisting}[style=jugeo-python, caption={Health monitoring for the advisor subsystem.}, label=lst:health]
from jugeo.kernel.health import (
    HealthStatus,
    HealthDimension,
    HealthIndicator,
    HealthReport,
)
from jugeo.evidence.channels import (
    CopilotChannel,
    ChannelConfiguration,
    EvidenceChannel,
    ChannelMonitor,
)

# Set up the Copilot channel and monitor
config = ChannelConfiguration.default_for(EvidenceChannel.COPILOT)
copilot = CopilotChannel(config=config)
monitor = ChannelMonitor()

# Check channel health
channel_healthy = copilot.health_check()

# Build health indicators for each dimension
connectivity = HealthIndicator(
    dimension=HealthDimension.COPILOT_CONNECTIVITY,
    status=HealthStatus.HEALTHY if channel_healthy else HealthStatus.UNHEALTHY,
    message="Copilot channel connectivity check",
    value=1.0 if channel_healthy else 0.0,
    timestamp=0.0,
    details={"channel": "copilot", "model": "copilot-default"},
)

evidence_flow = HealthIndicator(
    dimension=HealthDimension.EVIDENCE_FLOW,
    status=HealthStatus.HEALTHY,
    message="Evidence production rate within bounds",
    value=0.95,
    timestamp=0.0,
    details={"rate": monitor.throughput()},
)

trust_consistency = HealthIndicator(
    dimension=HealthDimension.TRUST_CONSISTENCY,
    status=HealthStatus.HEALTHY,
    message="All evidence respects trust ceiling",
    value=1.0,
    timestamp=0.0,
    details={"ceiling": "proposal", "violations": 0},
)

# Build the health report
report = HealthReport(
    overall_status=HealthStatus.HEALTHY,
    indicators=(connectivity, evidence_flow, trust_consistency),
    generated_at=0.0,
    report_id="advisor-health-001",
    degradation_reasons=(),
    recovery_suggestions=(),
)

# Query the report
print(f"Overall status: {report.overall_status}")
print(f"Fully healthy: {report.is_fully_healthy()}")
print(f"Indicator count: {len(report.indicators)}")

# Filter by status
healthy_indicators = report.filter_by_status(HealthStatus.HEALTHY)
print(f"Healthy indicators: {len(healthy_indicators)}")

# Filter by dimension
copilot_indicators = report.filter_by_dimension(
    HealthDimension.COPILOT_CONNECTIVITY
)
print(f"Copilot indicators: {len(copilot_indicators)}")

# Check worst indicator
worst = report.worst_indicator
print(f"Worst indicator: {worst.dimension} = {worst.status}")

# Degraded scenario
degraded_indicator = HealthIndicator(
    dimension=HealthDimension.COPILOT_CONNECTIVITY,
    status=HealthStatus.DEGRADED,
    message="Copilot response latency elevated",
    value=0.6,
    timestamp=0.0,
    details={"avg_latency_ms": 3500, "threshold_ms": 2000},
)

degraded_report = HealthReport(
    overall_status=HealthStatus.DEGRADED,
    indicators=(degraded_indicator, evidence_flow, trust_consistency),
    generated_at=0.0,
    report_id="advisor-health-002",
    degradation_reasons=("Elevated Copilot latency",),
    recovery_suggestions=("Increase timeout or switch to faster model",),
)
print(f"Degraded report: {degraded_report.overall_status}")
print(f"Degradation reasons: {degraded_report.degradation_reasons}")
print(f"Recovery suggestions: {degraded_report.recovery_suggestions}")
\end{lstlisting}

% ─────────────────────────────────────────────────────────────────────
\section{Formal Properties}
\label{sec:formal}

We prove five formal properties of the heap and scope advisors.
All proofs are formalized in \LEAN{} in
\texttt{Paper80\_HeapScopeAdvisors.lean}.

\subsection{Heap Advisory Soundness}
\label{sec:formal:heap}

\begin{theorem}[Heap advisory soundness]
\label{thm:heap-sound}
If the $\HeapAdv$ produces an advice item $a$ of kind
``aliasing hazard'' and $a$ passes validation by the solver
channel, then the identified aliasing pattern is indeed
hazardous: there exists an execution that mutates through
one alias while the other alias is live.
\end{theorem}

\begin{proof}
Validation by the solver channel means the solver found a model
(execution trace) in which both aliases are live and one is mutated.
This model is a concrete witness to the hazard.  Since the solver
is sound, the model is valid, and the aliasing pattern is
genuinely hazardous.

Formally, let $a = (\texttt{heap}, (r_1, r_2, \texttt{hazardous}), \gamma, \Tcopilot)$
be the advice item.  Validation constructs the formula:
\[
  \exists \sigma.\; \sigma(r_1) = \sigma(r_2)
  \;\wedge\; \mathrm{mutated}(\sigma, r_1)
  \;\wedge\; \mathrm{live}(\sigma, r_2)
\]
The solver confirms satisfiability, producing witness $\sigma^*$.
By soundness of the solver, $\sigma^*$ is a valid execution,
confirming the hazard.
\end{proof}

\subsection{Scope Advisory Soundness}
\label{sec:formal:scope}

\begin{theorem}[Scope advisory soundness]
\label{thm:scope-sound}
If the $\ScopeAdv$ produces an advice item $a$ of kind
``accidental shadowing'' and $a$ passes validation, then there
exists a scope $S_1$ and inner scope $S_2$ where the shadowed
variable $x$ in $S_2$ has different semantics from $x$ in $S_1$.
\end{theorem}

\begin{proof}
Validation checks that the shadowing satisfies the ``accidental''
criterion: the uses of $x$ in $S_2$ are semantically distinct from
the uses in $S_1$ (\ie renaming $x$ in $S_2$ does not affect $S_1$
but changes the program behavior, indicating the shadowing was
unintentional).  The solver verifies this by checking that:
\[
  \mathrm{rename}(S_2, x, x') \text{ changes the program output.}
\]
If this check passes, the shadowing is accidental.
\end{proof}

\subsection{Composition Trust Preservation}
\label{sec:formal:composition}

\begin{theorem}[Composition trust preservation]
\label{thm:comp-trust}
For any finite set of advisors $\{A_1, \ldots, A_n\}$ all using
the $\CopChan$, their composition $A_1 \otimes \cdots \otimes A_n$
produces advice with trust exactly $\Tcopilot$:
\[
  \tau(A_1 \otimes \cdots \otimes A_n) = \Tcopilot.
\]
\end{theorem}

\begin{proof}
By Lemma~\ref{lem:adv-trust}, each $\tau(A_i) = \Tcopilot$.
The composition trust is the iterated meet:
$\Tcopilot \tjoin \Tcopilot \tjoin \cdots \tjoin \Tcopilot = \Tcopilot$
by idempotence of $\tjoin$.
\end{proof}

\begin{corollary}[Adding advisors does not escalate trust]
\label{cor:no-escalation}
Adding new advisors to a composition never increases the trust
level of the combined advice.
\end{corollary}

\begin{proof}
By Theorem~\ref{thm:comp-trust}, the trust remains $\Tcopilot$
regardless of the number of advisors.
\end{proof}

\subsection{Advisory Completeness}
\label{sec:formal:completeness}

\begin{definition}[Advisory oracle]
\label{def:adv-oracle}
An \emph{advisory oracle} $\mathcal{O}_A$ for domain $d$ is a
function that, given any context $(c, \varphi, \Gamma)$, returns
the complete set of genuine issues in domain $d$.
\end{definition}

\begin{theorem}[Completeness relative to oracle]
\label{thm:completeness}
The $\HeapAdv$ is \emph{complete} relative to $\mathcal{O}_H$
if for every genuine heap issue $h \in \mathcal{O}_H(c, \varphi, \Gamma)$,
there exists an advice item $a \in \HeapAdv(c, \varphi, \Gamma)$
such that $a$ corresponds to $h$.

The $\HeapAdv$ using an actual LLM is not oracle-complete in
general.  However, for issues within the LLM's training
distribution, empirical completeness is typically high.
\end{theorem}

\begin{proof}
This is a conditional statement.  If the LLM detects all genuine
issues (as an oracle would), then the advisor is complete by
construction, since every detected issue generates an advice item.
The incompleteness gap arises from the LLM's imperfect detection.

Formally, let $D \subseteq \mathcal{O}_H(c, \varphi, \Gamma)$ be
the set of issues detected by the LLM.  Then the advisor produces
advice items for exactly $D$.  The completeness ratio is
$|D|/|\mathcal{O}_H|$.  Perfect completeness ($|D| = |\mathcal{O}_H|$)
requires oracle-level detection.
\end{proof}

\subsection{Health Monitoring Correctness}
\label{sec:formal:health}

\begin{theorem}[Health monitoring correctness]
\label{thm:health}
The health report produced by the advisor monitoring subsystem
accurately reflects the advisor's operational state:
\begin{enumerate}
\item If all indicators report $\mathsf{HEALTHY}$, the advisor
  is operational and producing evidence within bounds.
\item If any indicator reports $\mathsf{DEGRADED}$, the advisor
  is operational but with reduced quality or throughput.
\item If any indicator reports $\mathsf{UNHEALTHY}$, the advisor
  is non-functional for the affected dimension.
\end{enumerate}
\end{theorem}

\begin{proof}
Each health indicator is constructed from direct measurement of
the advisor's operational parameters:
\begin{itemize}
\item $\mathsf{COPILOT\_CONNECTIVITY}$ is measured by the
  $\CopChan.\texttt{health\_check}()$ method, which tests LLM
  responsiveness.
\item $\mathsf{EVIDENCE\_FLOW}$ is measured by the
  $\ChanMon.\texttt{throughput}()$ method.
\item $\mathsf{TRUST\_CONSISTENCY}$ is measured by checking
  that all produced evidence has trust $\leq \Tcopilot$.
\end{itemize}
Each measurement directly corresponds to the operational state
described in the theorem.  The status assignment (HEALTHY if
value $\geq 0.8$, DEGRADED if $0.4 \leq \mathrm{value} < 0.8$,
UNHEALTHY if $\mathrm{value} < 0.4$) maps measurements to
states.
\end{proof}

% ─────────────────────────────────────────────────────────────────────
\section{Discussion}
\label{sec:discussion}

\paragraph{Advisor effectiveness.}
The heap and scope advisors are most effective when the code
exhibits patterns that the LLM has seen during training.  For
idiomatic Python, Java, or TypeScript code, the LLM can reliably
identify common aliasing and shadowing patterns.  For unusual
coding patterns or domain-specific idioms, the advisor's
effectiveness may be reduced.

\paragraph{False positive management.}
LLM-generated advice may include false positives (flagged issues
that are not genuine).  The validation step filters these out,
but at the cost of solver invocations.  A high false positive
rate degrades overall system performance by wasting solver budget.

\paragraph{Compositional analysis.}
The composition of heap and scope advisors provides broader
coverage than either alone.  Cross-advisor interactions
(Section~\ref{sec:composition}) identify complex bugs that
span both heap and scope domains.

\paragraph{Scalability.}
The advisor pattern scales well with codebase size because the
LLM processes one context at a time.  The bottleneck is the
solver's ability to validate advice, not the LLM's ability to
generate it.

% ─────────────────────────────────────────────────────────────────────
\section{Related Work}
\label{sec:related}

\paragraph{Heap analysis.}
Static heap analysis has a long history, from separation
logic~\cite{reynolds2002separation} to shape
analysis~\cite{sagiv2002parametric}.  Our approach complements
these by using LLM hints to guide the analysis.

\paragraph{Scope analysis.}
Scope analysis in programming languages is well
understood~\cite{aho2006compilers}.  The novelty is using
LLM inference to classify shadowing as intentional or accidental.

\paragraph{LLM-assisted static analysis.}
Recent work uses LLMs to enhance static
analysis~\cite{chen2021codex,li2022competition}.
We formalize this enhancement within the trust architecture.

\paragraph{Advisor patterns.}
The advisor pattern in \JuGeo{} extends the concept of analysis
plugins found in IDE-based verification tools~\cite{leino2010dafny}.

% ─────────────────────────────────────────────────────────────────────
\section{Conclusion}
\label{sec:conclusion}

We have introduced the $\HeapAdv$ and $\ScopeAdv$, two
Copilot-guided analysis advisors within the \JuGeo{} framework.
The key results are:

\begin{enumerate}
\item \textbf{Heap advisory soundness} (Theorem~\ref{thm:heap-sound}).
\item \textbf{Scope advisory soundness} (Theorem~\ref{thm:scope-sound}).
\item \textbf{Composition trust preservation}
  (Theorem~\ref{thm:comp-trust}).
\item \textbf{Completeness relative to oracle}
  (Theorem~\ref{thm:completeness}).
\item \textbf{Health monitoring correctness}
  (Theorem~\ref{thm:health}).
\end{enumerate}

The advisor pattern provides a clean separation between issue
identification (LLM) and issue verification (solver/runtime).

% ─────────────────────────────────────────────────────────────────────
\begin{thebibliography}{25}

\bibitem{young-jugeo}
H.~Young.
\newblock Judgment Geometry: Semantic sites, descent, and the Copilot channel.
\newblock JuGeo Paper~0, 2024.

\bibitem{young-channels}
H.~Young.
\newblock The Copilot evidence channel.
\newblock JuGeo Paper~88, 2024.

\bibitem{young-advisors}
H.~Young.
\newblock The advisor architecture.
\newblock JuGeo Paper~90, 2024.

\bibitem{young-trust}
H.~Young.
\newblock The trust algebra.
\newblock JuGeo Paper~4, 2024.

\bibitem{young-heap}
H.~Young.
\newblock Heap aliasing in the judgment framework.
\newblock JuGeo Paper~18, 2024.

\bibitem{reynolds2002separation}
J.~C.~Reynolds.
\newblock Separation logic: A logic for shared mutable data structures.
\newblock In \emph{Proc.\ LICS}, 2002.

\bibitem{sagiv2002parametric}
S.~Sagiv, T.~W.~Reps, and R.~Wilhelm.
\newblock Parametric shape analysis via 3-valued logic.
\newblock \emph{TOPLAS}, 24(3):217--298, 2002.

\bibitem{aho2006compilers}
A.~V.~Aho, M.~S.~Lam, R.~Sethi, and J.~D.~Ullman.
\newblock \emph{Compilers: Principles, Techniques, and Tools}.
\newblock Pearson, 2nd edition, 2006.

\bibitem{chen2021codex}
M.~Chen, J.~Tworek, H.~Jun, et~al.
\newblock Evaluating large language models trained on code.
\newblock \emph{arXiv:2107.03374}, 2021.

\bibitem{li2022competition}
Y.~Li, D.~Choi, J.~Chung, et~al.
\newblock Competition-level code generation with {AlphaCode}.
\newblock \emph{Science}, 378(6624):1092--1097, 2022.

\bibitem{leino2010dafny}
K.~R.~M.~Leino.
\newblock {Dafny}: An automatic program verifier.
\newblock In \emph{Proc.\ LPAR}, 2010.

\bibitem{hoare1969axiomatic}
C.~A.~R.~Hoare.
\newblock An axiomatic basis for computer programming.
\newblock \emph{CACM}, 12(10):576--580, 1969.

\bibitem{demoura2008z3}
L.~de~Moura and N.~Bj{\o}rner.
\newblock {Z3}: An efficient {SMT} solver.
\newblock In \emph{Proc.\ TACAS}, 2008.

\bibitem{de2015lean}
L.~de~Moura, S.~Kong, J.~Avigad, F.~van~Doorn, and J.~von~Raumer.
\newblock The {Lean} theorem prover.
\newblock In \emph{Proc.\ CADE}, 2015.

\bibitem{ouyang2022training}
L.~Ouyang, J.~Wu, X.~Jiang, et~al.
\newblock Training language models to follow instructions.
\newblock In \emph{Proc.\ NeurIPS}, 2022.

\bibitem{dijkstra1975guarded}
E.~W.~Dijkstra.
\newblock Guarded commands and formal derivation.
\newblock \emph{CACM}, 18(8):453--457, 1975.

\end{thebibliography}

\appendix
\section{Additional Proofs}
\label{app:proofs}

\subsection{Advisor Category Structure}
\label{app:category}

\begin{definition}[Category of advisors]
\label{def:adv-cat}
The \emph{category of advisors} $\mathbf{Adv}$ has:
\begin{itemize}[leftmargin=2em]
\item \textbf{Objects:} Verification contexts $(c, \varphi, \Gamma)$.
\item \textbf{Morphisms:} Advisor functions $A : (c_1, \varphi_1, \Gamma_1) \to \mathcal{P}(\text{advice})$.
\item \textbf{Composition:} Union of advice sets with trust meet.
\item \textbf{Identity:} The trivial advisor that produces no advice.
\end{itemize}
\end{definition}

\begin{lemma}[Advisor composition is associative]
\label{lem:adv-assoc}
$(A_1 \otimes A_2) \otimes A_3 = A_1 \otimes (A_2 \otimes A_3)$.
\end{lemma}

\begin{proof}
Both sides produce the union of advice sets from $A_1$, $A_2$,
and $A_3$.  Union is associative.  The trust meet is also
associative by the lattice property.
\end{proof}

\begin{lemma}[Advisor composition is commutative]
\label{lem:adv-comm}
$A_1 \otimes A_2 = A_2 \otimes A_1$.
\end{lemma}

\begin{proof}
Union and meet are commutative.
\end{proof}

\begin{lemma}[Identity advisor is neutral]
\label{lem:adv-id}
$A \otimes \mathrm{id} = A$ for any advisor $A$.
\end{lemma}

\begin{proof}
The identity advisor produces no advice ($\emptyset$).
$A_{\text{advice}} \cup \emptyset = A_{\text{advice}}$.
The trust is $\tau(A) \tjoin \Tproof = \tau(A)$ (since $\Tproof$
is the top element and $\min(\tau, \Tproof) = \tau$).
\end{proof}

\subsection{Health Status Lattice}
\label{app:health-lattice}

\begin{lemma}[Health status is a lattice]
\label{lem:health-lattice}
The health statuses
$\mathsf{UNHEALTHY} < \mathsf{DEGRADED} < \mathsf{UNKNOWN}
< \mathsf{RECOVERING} < \mathsf{HEALTHY}$
form a totally ordered lattice.
\end{lemma}

\begin{proof}
A finite totally ordered set is a lattice with meet $= \min$
and join $= \max$.
\end{proof}

\begin{lemma}[Health report status is the meet of indicators]
\label{lem:health-meet}
The overall status of a health report with indicators
$h_1, \ldots, h_n$ is $\min(h_1.\mathrm{status}, \ldots, h_n.\mathrm{status})$.
\end{lemma}

\begin{proof}
The report takes the worst (minimum) status among its indicators,
which is the meet in the health status lattice.
\end{proof}

\begin{lemma}[Adding indicators can only degrade status]
\label{lem:health-degrade}
If a new indicator $h_{n+1}$ is added to a report, the overall
status can only stay the same or decrease:
$\min(h_1, \ldots, h_n, h_{n+1}) \leq \min(h_1, \ldots, h_n)$.
\end{lemma}

\begin{proof}
$\min(S \cup \{x\}) \leq \min(S)$ for any set $S$ and element $x$.
\end{proof}

\subsection{Heap Advisor Lattice Properties}
\label{app:heap-lattice}

\begin{lemma}[Heap state forms a partial order]
\label{lem:heap-po}
The heap state preorder $(\mathcal{H}, \sqsubseteq)$, where
$h \sqsubseteq h'$ iff $h'$ extends $h$ with additional
allocation records, is a partial order.
\end{lemma}

\begin{proof}
\emph{Reflexivity:} $h \sqsubseteq h$ since every heap extends
itself trivially.
\emph{Antisymmetry:} if $h \sqsubseteq h'$ and $h' \sqsubseteq h$,
then $h$ and $h'$ have the same allocation records, so $h = h'$.
\emph{Transitivity:} if $h \sqsubseteq h'$ and $h' \sqsubseteq h''$,
then $h''$ extends $h'$ which extends $h$, so $h \sqsubseteq h''$.
\end{proof}

\begin{lemma}[Heap advisor monotonicity]
\label{lem:heap-mono}
If $h \sqsubseteq h'$, then the set of heap advisor suggestions
for $h$ is a subset of the suggestions for $h'$:
$\mathrm{suggest}(h) \subseteq \mathrm{suggest}(h')$.
\end{lemma}

\begin{proof}
A larger heap state provides more information for the advisor.
Every suggestion valid for $h$ remains valid for $h'$ since $h'$
contains all records of $h$.  Additional records in $h'$ may
generate additional suggestions but cannot invalidate existing ones.
\end{proof}

\begin{definition}[Scope nesting depth]
\label{def:scope-depth}
The \emph{nesting depth} of a scope $s$ is the length of the
longest chain $s_0 \subsetneq s_1 \subsetneq \cdots \subsetneq s$
in the scope containment order.
\end{definition}

\begin{lemma}[Scope nesting is well-founded]
\label{lem:scope-wf}
For a finite program, the scope nesting depth is bounded by the
maximum syntactic nesting depth of the program.
\end{lemma}

\begin{proof}
Each scope corresponds to a syntactic block (function, class, loop,
conditional).  The nesting depth of scopes equals the nesting
depth of the corresponding blocks.  Since the program is finite,
the nesting depth is bounded.
\end{proof}

\begin{lemma}[Combined advisor trust ceiling]
\label{lem:combined-ceil}
When the heap advisor and scope advisor are composed, the trust
ceiling of the composition is $\min(\Tcopilot, \Tcopilot) = \Tcopilot$.
\end{lemma}

\begin{proof}
Both advisors have trust ceiling $\Tcopilot$.  Sequential or
parallel composition takes the minimum, which is $\Tcopilot$.
\end{proof}

\begin{proposition}[Heap-scope advisory completeness]
\label{prop:hs-complete}
For any heap state $h$ and scope $s$, the combined heap-scope
advisor produces a non-empty suggestion list if either $h$ contains
at least one allocation or $s$ has at least one binding.
\end{proposition}

\begin{proof}
If $h$ has at least one allocation, the heap advisor generates
at least a ``check allocation validity'' suggestion.  If $s$ has
at least one binding, the scope advisor generates at least a
``check binding visibility'' suggestion.  In either case, the
combined suggestion list is non-empty.
\end{proof}

\begin{remark}[Extension to interprocedural analysis]
The heap and scope advisors operate intraprocedurally.
Extending them to interprocedural analysis requires tracking
heap effects across call boundaries, which introduces additional
complexity.  We leave this extension to future work.
\end{remark}

\end{document}
